\section{非线性数据结构章正文案例推导补强}

这一节补齐本章正文出现过的 LeetCode 案例推导。阅读顺序统一按“题目说明、样例入口、暴力基线、暴力过程、重复工作、优化条件、相邻状态、状态复用、指针和字段、代码落点、正确性、边界实验、复杂度变化、迁移追问”展开；目的不是新增题量，而是把正文中的案例都讲成可复述、可证明、可迁移的解题链。

\subsection{LeetCode 1 两数之和}

\textbf{题目说明。} LeetCode 1 两数之和 的题面入口要先还原成具体任务：当前元素只需要寻找唯一补数，历史值到下标的映射就是最小状态。

\textbf{样例入口。} 拿数组 [2, 7]、target=9 手算：站在 2 时历史为空，不能配对；站在 7 时只需要问历史里有没有 2。再试 [3, 3]、target=6：如果先把当前 3 插入再查，就会把同一个下标用两次。

\textbf{暴力基线。} 暴力两层循环枚举 i 和 j，检查 nums[i]+nums[j] 是否等于 target，找到后返回两个下标。

\textbf{暴力过程。} 以 nums=[2, 7, 11, 15]、target=9 为例，暴力先试 2+7 命中；换成 [3, 2, 4]、target=6 时，会试 3+2 失败，再试 2+4 命中。

\textbf{重复工作。} 题面模型：当前元素只需要寻找唯一补数，历史值到下标的映射就是最小状态。暴力版的慢点要落到本题：浪费在每个当前位置都回头扫描历史。站在 nums[i] 时，真正需要的不是所有历史值，而是历史里有没有 target-nums[i]。后面的优化只消掉这类重复工作，不能改变答案集合。

\textbf{优化条件。} 本题的关键观察是：先查再插入，避免同一个下标被用两次；若数组有序才考虑双指针。这句话必须能翻译成可维护状态；如果题目条件改变后观察不再成立，就不能继续套原来的写法。

\textbf{相邻状态。} 规律不是“用哈希表”四个字，而是每个当前位置只和它左侧历史配对，代码必须先查补数再插入当前值。把这个特例继续拆开看：每个合法答案都要还原成当前状态和某个历史状态的配对；哈希表的 key 负责定位历史状态，value 负责表达次数、最早位置或存在性。先查还是先写，必须由是否允许当前前缀和自己配对来决定。

\textbf{状态复用。} 哈希表保存已经扫描过的值到下标的映射。每到一个新数，先查补数是否存在；不存在再把当前数放入历史。手算时只改被新输入影响的那一小部分，而不是回头重算完整候选。

\textbf{指针和字段。} 状态字段要能说清：i 是当前下标，need=target-nums[i] 是补数，seen[value] 是该值在左侧历史中出现的下标。手算时按这个协议跟踪：nums=[3, 3]、target=6 时，第一个 3 先查不到补数，再插入；第二个 3 查到历史 3，返回两个不同下标。

\textbf{代码落点。} 关键顺序是先 find(need)，再 seen[nums[i]]=i。若先插入当前值，target=2*nums[i] 时可能把同一位置用两次。关键代码行必须能逐句翻译成“旧状态怎样变成新状态”，否则就只是模板。

\textbf{正确性。} 每个答案对都有一个较右下标；扫描到较右下标时，较左下标已经在 seen 中，所以一定能被查到。

\textbf{边界实验。} 先用重复数字、目标为偶数且补数等于当前值、没有答案时返回空结果做反例检查。实验时把重复数字、目标为偶数且补数等于当前值、没有答案时返回空结果加入对拍集合，用平方暴力和优化版比较。失败时先看初始历史状态、查询更新顺序和哈希表 value 类型。

\textbf{复杂度变化。} 暴力 O(\ensuremath{n^{2}})，哈希表扫描 O(n) 均摊时间、O(n) 空间。

\textbf{迁移追问。} 如果题目改成“若改成返回所有不重复数对，要先排序或用频次表处理重复”，先判断原状态是否仍然覆盖新答案集合，再决定是微调边界、改 value 语义，还是换算法。

\subsection{LeetCode 560 和为 K 的子数组}

\textbf{题目说明。} LeetCode 560 和为 K 的子数组给定整数数组 nums 和整数 k，统计和等于 k 的连续子数组个数；数组中可以有负数。

\textbf{样例入口。} 样例 nums=[1, 1, 1]，k=2，输出 2；两个合法连续子数组分别是前两个 1 和后两个 1。

\textbf{暴力基线。} 暴力固定 left，向右累计 sum；每当 sum==k 就计数。

\textbf{暴力过程。} 以 nums=[1, 1, 1]、k=2 为例，从 left=0 可找到 [1, 1]，从 left=1 又找到后两个 1。

\textbf{重复工作。} 题面模型：当前前缀和减去目标值，就是需要匹配的历史前缀。暴力版的慢点要落到本题：浪费在每个右端都回头枚举所有 left。区间和等于 k 可以改写成 prefix[right]-prefix[left]=k。后面的优化只消掉这类重复工作，不能改变答案集合。

\textbf{优化条件。} 本题的关键观察是：哈希表保存前缀和出现次数，先查询再更新。这句话必须能翻译成可维护状态；如果题目条件改变后观察不再成立，就不能继续套原来的写法。

\textbf{相邻状态。} 规律是哈希表保存前缀和频次，且初始要放 0:1。把这个特例继续拆开看：每个合法答案都要还原成当前状态和某个历史状态的配对；哈希表的 key 负责定位历史状态，value 负责表达次数、最早位置或存在性。先查还是先写，必须由是否允许当前前缀和自己配对来决定。

\textbf{状态复用。} 扫描前缀和 prefix。哈希表 count[old\_prefix] 保存历史前缀出现次数；当前位置需要查 prefix-k 出现过几次。手算时只改被新输入影响的那一小部分，而不是回头重算完整候选。

\textbf{指针和字段。} 状态字段要能说清：prefix 是当前前缀和，need=prefix-k 是所需历史前缀，count 保存历史前缀频次，answer 累加匹配次数。手算时按这个协议跟踪：[1, 1, 1] 中扫描到第二个元素时 prefix=2，需要历史 0，count[0]=1，所以找到第一个长度 2 子数组。

\textbf{代码落点。} 关键是先 answer += count[prefix-k]，再 count[prefix]++。初始化 count[0]=1，表示从下标 0 开始的区间。关键代码行必须能逐句翻译成“旧状态怎样变成新状态”，否则就只是模板。

\textbf{正确性。} 任意和为 k 的子数组都对应一对前缀差 k；扫描到右端时，所有合法左端前缀都已在表里。

\textbf{边界实验。} 先用负数、目标为零、从下标零开始的区间、重复前缀做反例检查。实验时把负数、目标为零、从下标零开始的区间、重复前缀加入对拍集合，用平方暴力和优化版比较。失败时先看初始历史状态、查询更新顺序和哈希表 value 类型。

\textbf{复杂度变化。} 暴力 O(\ensuremath{n^{2}})，前缀哈希 O(n) 均摊时间、O(n) 空间。

\textbf{迁移追问。} 如果题目改成“若要求最长区间，哈希表应保存最早位置而不是次数”，先判断原状态是否仍然覆盖新答案集合，再决定是微调边界、改 value 语义，还是换算法。

\subsection{LeetCode 215 数组中的第 K 个最大元素}

\textbf{题目说明。} LeetCode 215 数组中的第 K 个最大元素 的题面入口要先还原成具体任务：只关心第 K 大，不需要完整排序。

\textbf{样例入口。} 拿 [3, 2, 1, 5, 6, 4]、k=2 手算，固定大小为 2 的小顶堆最终保存最大的两个数，堆顶就是第 2 大。

\textbf{暴力基线。} 暴力排序整个数组，然后返回降序第 k 个或升序第 n-k 个。

\textbf{暴力过程。} 以 [3, 2, 1, 5, 6, 4]、k=2 为例，完整排序后第 2 大是 5；但为了一个第 K 大，不需要知道所有元素的完整顺序。

\textbf{重复工作。} 题面模型：只关心第 K 大，不需要完整排序。暴力版的慢点要落到本题：浪费在完整排序。我们只需要维护当前最大的 K 个元素，其中最小的那个就是第 K 大边界。后面的优化只消掉这类重复工作，不能改变答案集合。

\textbf{优化条件。} 本题的关键观察是：大小为 K 的小顶堆保存当前前 K 大边界；快速选择可做到平均线性。这句话必须能翻译成可维护状态；如果题目条件改变后观察不再成立，就不能继续套原来的写法。

\textbf{相邻状态。} 规律是堆里只保留当前前 k 大候选，比堆顶小的数不可能进入答案。把这个特例继续拆开看：堆顶必须正好代表下一步最需要处理的候选；若堆里可能有过期对象，读取堆顶前要先清理。堆只解决动态极值，不自动证明被弹出或被丢弃的元素无用。

\textbf{状态复用。} 大小为 K 的小顶堆保存当前前 K 大。新数大于堆顶时，弹出堆顶并加入新数；小于等于堆顶则不可能进入前 K。手算时只改被新输入影响的那一小部分，而不是回头重算完整候选。

\textbf{指针和字段。} 状态字段要能说清：heap.top 是当前前 K 大中的最小值，size 保持不超过 K，num 是当前扫描元素。手算时按这个协议跟踪：扫描样例时，堆最终保存 5 和 6，堆顶 5 就是第 2 大。

\textbf{代码落点。} 关键是使用小顶堆，C++ priority\_queue 默认大顶堆，需要 greater<int>。若选择快速选择，要处理随机 pivot 和重复值。关键代码行必须能逐句翻译成“旧状态怎样变成新状态”，否则就只是模板。

\textbf{正确性。} 堆中始终保存已扫描元素的前 K 大；若新元素不超过堆顶，它最多排在第 K 大之后，可以安全丢弃。

\textbf{边界实验。} 先用K 为一、K 为 n、重复值、随机 pivot做反例检查。实验时覆盖K 为一、K 为 n、重复值、随机 pivot，记录堆顶是否过期、比较器是否让正确候选在堆顶、K 或窗口大小为边界值时是否稳定。

\textbf{复杂度变化。} 排序 O(n log n)，小顶堆 O(n log k) 时间、O(k) 空间；快速选择平均 O(n)。

\textbf{迁移追问。} 如果题目改成“若数据流持续到来，堆方案比一次性快速选择更自然”，先判断原状态是否仍然覆盖新答案集合，再决定是微调边界、改 value 语义，还是换算法。

\subsection{LeetCode 98 验证二叉搜索树}

\textbf{题目说明。} LeetCode 98 验证二叉搜索树 的题面入口要先还原成具体任务：BST 约束是全局上下界，不是只比较父子节点。

\textbf{样例入口。} 拿 [5, 1, 4, null, null, 3, 6] 手算，节点 3 小于根 5 却在右子树里，只比较父子 4 和 3 会误判。

\textbf{暴力基线。} 错误暴力常只检查当前节点大于左孩子、小于右孩子；正确慢版本要对每个节点扫描整棵左子树最大值和右子树最小值。

\textbf{暴力过程。} 以 [5, 1, 4, null, null, 3, 6] 为例，节点 4 的左孩子 3 小于 4、右孩子 6 大于 4，看局部没问题；但 3 在根 5 的右子树里，小于 5，所以整棵树不是 BST。

\textbf{重复工作。} 题面模型：BST 约束是全局上下界，不是只比较父子节点。暴力版的慢点要落到本题：浪费在每个节点重复扫描子树求上下界。真正需要传递的是祖先给当前子树限定的取值范围。后面的优化只消掉这类重复工作，不能改变答案集合。

\textbf{优化条件。} 本题的关键观察是：中序严格递增或显式传递上下界都可以；中序版本要保存前一个访问值。这句话必须能翻译成可维护状态；如果题目条件改变后观察不再成立，就不能继续套原来的写法。

\textbf{相邻状态。} 规律是 BST 约束来自整条祖先上下界，或中序序列必须严格递增；重复值和 int 边界要按题意单独处理。把这个特例继续拆开看：节点向父亲返回的量和全局答案通常不是同一个量。先写叶子或空节点的返回值，再写父节点如何合并左右孩子，递归式才有边界基础。

\textbf{状态复用。} 自顶向下传入 lower 和 upper，当前节点值必须满足 lower<val<upper；去左子树时 upper 变成当前值，去右子树时 lower 变成当前值。手算时只改被新输入影响的那一小部分，而不是回头重算完整候选。

\textbf{指针和字段。} 状态字段要能说清：node 是当前节点，lower/upper 是祖先约束；空节点返回合法。手算时按这个协议跟踪：到样例中右子树节点 3 时，它的 lower 应该是 5、upper 是 4 或来自路径约束，3 不满足根给的下界。

\textbf{代码落点。} 关键是用 long long 或 optional 边界，避免节点值等于 int 最小/最大时哨兵冲突；重复值是否允许按题意处理，本题严格不允许。关键代码行必须能逐句翻译成“旧状态怎样变成新状态”，否则就只是模板。

\textbf{正确性。} BST 要求左子树所有值小于根、右子树所有值大于根；上下界把所有祖先约束都带到当前节点，覆盖了非局部违规。

\textbf{边界实验。} 先用重复值、int 最小最大值、局部合法但全局非法做反例检查。实验时覆盖重复值、int 最小最大值、局部合法但全局非法，尤其关注空节点、单节点、链式树和全负值。递归版本和显式栈版本可以互相对拍。

\textbf{复杂度变化。} 扫描子树慢版本最坏 O(\ensuremath{n^{2}})，上下界 DFS/BFS 每节点一次 O(n)，空间 O(h) 或显式栈 O(h)。

\textbf{迁移追问。} 如果题目改成“若允许重复值，需要明确重复放左还是放右”，先判断原状态是否仍然覆盖新答案集合，再决定是微调边界、改 value 语义，还是换算法。

\subsection{LeetCode 207 课程表}

\textbf{题目说明。} LeetCode 207 课程表 的题面入口要先还原成具体任务：课程是节点，先修关系是有向边，问题是判断是否有环。

\textbf{样例入口。} 拿课程 0->1 手算，入度为 0 的课程 0 先入队，学完后 1 的入度降为 0；若有 0->1 和 1->0，队列一开始为空。

\textbf{暴力基线。} 暴力可以不断扫描所有课程，找当前先修都已满足的课程；若一轮没有新课程可学，就判断有环。

\textbf{暴力过程。} 以 0->1 为例，课程 0 入度为 0，先学 0，随后 1 的入度变 0。若有 0->1 和 1->0，一开始没有入度为 0 的课。

\textbf{重复工作。} 题面模型：课程是节点，先修关系是有向边，问题是判断是否有环。暴力版的慢点要落到本题：浪费在每轮都重新检查所有课程的先修。入度正好表示还剩多少未满足依赖，学完一门课只影响它指向的后继。后面的优化只消掉这类重复工作，不能改变答案集合。

\textbf{优化条件。} 本题的关键观察是：入度拓扑排序不断学习当前无依赖课程。这句话必须能翻译成可维护状态；如果题目条件改变后观察不再成立，就不能继续套原来的写法。

\textbf{相邻状态。} 规律是拓扑排序用入度为 0 的节点推进，处理数量小于课程数说明有环。把这个特例继续拆开看：状态节点不一定等于原始位置，还可能包含钥匙、剩余资源、时间或访问集合。visited 只能合并未来能力完全相同的状态，否则会把合法路径剪掉。

\textbf{状态复用。} 构建邻接表和 indegree。把所有入度为 0 的课程入队，弹出课程时让后继入度减一，减到 0 就入队。手算时只改被新输入影响的那一小部分，而不是回头重算完整候选。

\textbf{指针和字段。} 状态字段要能说清：indegree[v] 是课程 v 尚未满足的先修数量，queue 保存当前可学习课程，visited\_count 是已处理数量。手算时按这个协议跟踪：0 出队后，遍历它的后继 1，indegree[1] 从 1 变 0，1 入队；最后处理数量等于课程数，说明无环。

\textbf{代码落点。} 关键是边方向要统一：prerequisite [a, b] 表示 b -> a。处理数量小于 numCourses 时返回 false。关键代码行必须能逐句翻译成“旧状态怎样变成新状态”，否则就只是模板。

\textbf{正确性。} 入度为 0 的节点没有未满足前置条件，可以安全学习；有向环中的节点永远不会全部降到入度 0。

\textbf{边界实验。} 先用边方向、孤立课程、环、自依赖、重复边做反例检查。实验时覆盖边方向、孤立课程、环、自依赖、重复边，并打印入队时的完整状态。若同一位置但资源不同的状态被合并，很多图题会出现隐蔽错误。

\textbf{复杂度变化。} 反复扫描最坏 O(VE)，拓扑排序 O(V+E) 时间和空间。

\textbf{迁移追问。} 如果题目改成“DFS 三色标记也能判环，但要明确访问状态”，先判断原状态是否仍然覆盖新答案集合，再决定是微调边界、改 value 语义，还是换算法。

\subsection{LeetCode 200 岛屿数量}

\textbf{题目说明。} LeetCode 200 岛屿数量 的题面入口要先还原成具体任务：陆地格子是节点，四方向相邻是边，岛屿是连通块。

\textbf{样例入口。} 拿一个 2x2 小岛手算，从一个陆地出发能沿上下左右标记同一块岛。若不标记访问，会重复计数；若对角线也连通，会误合并。

\textbf{暴力基线。} 暴力可以对每个陆地格子单独搜索它所属的连通块，再用集合去重；这样同一块岛会被重复搜索很多次。

\textbf{暴力过程。} 以 2x2 全是陆地为例，从左上搜索会访问四个格子；如果之后从右上又重新搜索，就把同一座岛重复数了一遍。

\textbf{重复工作。} 题面模型：陆地格子是节点，四方向相邻是边，岛屿是连通块。暴力版的慢点要落到本题：浪费在已确认属于某个岛的陆地没有被标记。每块岛只需要在第一次遇到时完整淹没或标记。后面的优化只消掉这类重复工作，不能改变答案集合。

\textbf{优化条件。} 本题的关键观察是：扫描遇到未访问陆地就启动搜索并标记整块。这句话必须能翻译成可维护状态；如果题目条件改变后观察不再成立，就不能继续套原来的写法。

\textbf{相邻状态。} 规律是每发现一个未访问陆地，答案加一，并用 DFS/BFS 标记与它四连通的全部陆地。把这个特例继续拆开看：状态节点不一定等于原始位置，还可能包含钥匙、剩余资源、时间或访问集合。visited 只能合并未来能力完全相同的状态，否则会把合法路径剪掉。

\textbf{状态复用。} 扫描网格。遇到未访问陆地，答案加一，然后 BFS/DFS 把与它四方向连通的全部陆地标记为已访问。手算时只改被新输入影响的那一小部分，而不是回头重算完整候选。

\textbf{指针和字段。} 状态字段要能说清：row/col 是当前格子，visited 或原地改写表示已经归属某座岛，directions 是四方向邻居。手算时按这个协议跟踪：从左上陆地入队，依次扩展右边和下边陆地，最终整块连通区域都标记；扫描到这些格子时直接跳过。

\textbf{代码落点。} 关键是只看上下左右，不看对角线。边界判断集中写，访问标记在入队时设置，避免重复入队。关键代码行必须能逐句翻译成“旧状态怎样变成新状态”，否则就只是模板。

\textbf{正确性。} 一次搜索会访问且只访问与起点四连通的全部陆地；外层每次启动搜索都对应一块此前未计数的新岛。

\textbf{边界实验。} 先用空网格、全水、全陆、斜对角不连通、是否允许修改输入做反例检查。实验时覆盖空网格、全水、全陆、斜对角不连通、是否允许修改输入，并打印入队时的完整状态。若同一位置但资源不同的状态被合并，很多图题会出现隐蔽错误。

\textbf{复杂度变化。} 重复搜索最坏 O(\ensuremath{(mn)^{2}})，标记搜索 O(mn) 时间，空间 O(mn) 或递归/队列大小。

\textbf{迁移追问。} 如果题目改成“也可用并查集合并相邻陆地，适合动态岛屿扩展”，先判断原状态是否仍然覆盖新答案集合，再决定是微调边界、改 value 语义，还是换算法。

\subsection{LeetCode 684 冗余连接}

\textbf{题目说明。} LeetCode 684 冗余连接 的题面入口要先还原成具体任务：树加一条边会形成唯一环。

\textbf{样例入口。} 拿 edges=[[1, 2], [1, 3], [2, 3]] 手算，前两条边把 1、2、3 连成一棵树，第三条边连接的两个点已经同组，所以它就是冗余边。

\textbf{暴力基线。} 慢版本每次合并后用 DFS 或 BFS 重新判断连通性。它正确但重复遍历已有连接；当关系只会合并不会拆开时，可以压成集合代表。放到本题就是：先按“树加一条边会形成唯一环”完整试慢版本；样例里已经能看到“规律是无向图加边时，第一次 union 失败的边形成环”，所以优化前必须先能说清慢版本枚举了哪些候选。

\textbf{暴力过程。} 拿 edges=[[1, 2], [1, 3], [2, 3]] 手算，前两条边把 1、2、3 连成一棵树，第三条边连接的两个点已经同组，所以它就是冗余边。

\textbf{重复工作。} 题面模型：树加一条边会形成唯一环。暴力版的慢点要落到本题：慢版本会围绕“树加一条边会形成唯一环”反复搜索连通性。样例规律是“规律是无向图加边时，第一次 union 失败的边形成环”，说明关系只需要被压成集合代表；每次 union 失败或成功都要能翻译回题意。后面的优化只消掉这类重复工作，不能改变答案集合。

\textbf{优化条件。} 本题的关键观察是：按顺序加边，若一条边两端已连通，则它就是冗余边。这句话必须能翻译成可维护状态；如果题目条件改变后观察不再成立，就不能继续套原来的写法。

\textbf{相邻状态。} 规律是无向图加边时，第一次 union 失败的边形成环。把这个特例继续拆开看：union 只是在维护等价类，不是在保存具体路径。两个节点代表元相同只能说明连通或关系一致；如果题目还要距离、比例或奇偶性，必须把附加信息挂到根关系上。

\textbf{状态复用。} 把连通分量压成代表元；查询只比较代表，合并只发生在两个根之间。本题落点是：规律是无向图加边时，第一次 union 失败的边形成环。手算时只改被新输入影响的那一小部分，而不是回头重算完整候选。

\textbf{指针和字段。} 状态字段要能说清：parent、size 或 rank、find 返回的代表元、合并时的附加信息；带权或奇偶并查集还要维护到根的关系。手算时按这个协议跟踪：拿 edges=[[1, 2], [1, 3], [2, 3]] 手算，前两条边把 1、2、3 连成一棵树，第三条边连接的两个点已经同组，所以它就是冗余边。然后按协议复查：手算时写 parent 表和集合代表。每次 union 只改变根之间的关系，find 后代表相同才说明连通。

\textbf{代码落点。} find 路径压缩不能破坏附加权值；union 只能连接两个根；合并失败和重复边要保留题意解释。关键代码行必须能逐句翻译成“旧状态怎样变成新状态”，否则就只是模板。

\textbf{正确性。} 证明从等价关系开始：合并维护连通性的传递性，find 返回的代表元就是当前集合身份。

\textbf{边界实验。} 先用节点编号、最后一条成环边、重复边、并查集初始化大小做反例检查。实验时覆盖节点编号、最后一条成环边、重复边、并查集初始化大小，打印每个元素的代表元和集合大小。重复合并、已经连通的边和字符串编号最容易暴露实现问题。

\textbf{复杂度变化。} 暴力每次查询都搜索连通性代价高；并查集把合并和查询摊还到近似常数，整体接近线性。

\textbf{迁移追问。} 如果题目改成“有向冗余连接需要同时处理入度为二和环，复杂度更高”，先判断原状态是否仍然覆盖新答案集合，再决定是微调边界、改 value 语义，还是换算法。

\begin{principlebox}{本节复盘方式}
不要只看最终算法名。每道题复盘时先遮住优化版，口述暴力枚举对象；再指出相邻窗口、历史前缀、候选区间、搜索状态或子问题里哪一部分被重复处理；最后用一个边界样例验证状态更新顺序。能讲完这三步，才算真正掌握本章案例。
\end{principlebox}
